\section{Návrh používateľského rozhrania}

Používateľské rozhranie manipulátora slúži na zadávanie požadovanej polohy, orientácie alebo priamo kĺbových uhlov robota. 
Rozhranie je realizované pomocou LCD displeja 16x2 s I2C prevodníkom a joysticku s tlačidlami.

Systém beží ako samostatná FreeRTOS úloha \texttt{TaskUI}, ktorá zabezpečuje:
\begin{itemize}
    \item čítanie vstupu z joysticku,
    \item prepínanie medzi obrazovkami,
    \item prepínanie medzi módmi ovládania,
    \item zobrazovanie aktuálneho stavu robota,
    \item odosielanie požiadaviek do riadiacej úlohy.
\end{itemize}

\subsection{Hardvérové prvky}

Používateľské rozhranie využíva nasledovné hardvérové prvky:
\begin{itemize}
    \item \textbf{Joystick} -- slúži na zmenu hodnôt v dvoch osiach.
    \item \textbf{Tlačidlo A} -- prepína medzi jednotlivými obrazovkami.
    \item \textbf{Tlačidlo B} -- prepína medzi módmi v rámci obrazovky.
    \item \textbf{LCD 16x2} -- zobrazuje aktuálne hodnoty.
\end{itemize}

\subsection{Obrazovky}

Používateľské rozhranie obsahuje tri základné obrazovky:

\subsubsection{Obrazovka polohy XYZ}

Na tejto obrazovke sa nastavuje poloha koncového efektora v priestore.

\begin{lstlisting}[caption={Ukážka obrazovky XYZ}]
X:12.4   Y:8.1
Z:4.2    XY
\end{lstlisting}

V pravom dolnom rohu sa zobrazuje aktuálny mód:

\begin{itemize}
    \item \texttt{XY} -- joystick ovláda osi X a Y
    \item \texttt{YZ} -- joystick ovláda osi Y a Z
    \item \texttt{ZX} -- joystick ovláda osi Z a X
\end{itemize}

\subsubsection{Obrazovka orientácie RPY}

Na tejto obrazovke sa nastavuje orientácia efektora pomocou Roll, Pitch a Yaw uhlov.

\begin{lstlisting}[caption={Ukážka obrazovky RPY}]
R:15.0   P:8.2
Y:4.1    RP
\end{lstlisting}

Možné módy:

\begin{itemize}
    \item \texttt{RP} -- joystick ovláda Roll a Pitch
    \item \texttt{PY} -- joystick ovláda Pitch a Yaw
    \item \texttt{YR} -- joystick ovláda Yaw a Roll
\end{itemize}

\subsubsection{Obrazovka kĺbových uhlov}

Na tejto obrazovke sa nastavujú jednotlivé uhly servomotorov.

\begin{lstlisting}[caption={Ukážka obrazovky JOINTS}]
T1:20.0  T2:10.0
PAIR:12
\end{lstlisting}

Možné módy:

\begin{itemize}
    \item \texttt{PAIR:12} -- joystick ovláda $\theta_1$ a $\theta_2$
    \item \texttt{PAIR:34} -- joystick ovláda $\theta_3$ a $\theta_4$
    \item \texttt{PAIR:56} -- joystick ovláda $\theta_5$ a $\theta_6$
\end{itemize}

\subsection{Spracovanie joysticku}

Joystick poskytuje analógové hodnoty v rozsahu približne 0 až 1023. 
Stredová hodnota je približne:

\begin{equation}
CENTER = 520
\end{equation}

Kvôli odstráneniu šumu sa používa mŕtva zóna:

\begin{equation}
DEADZONE = 20
\end{equation}

Výpočet zmeny:

\begin{equation}
dx = (joyX - CENTER)\cdot k
\end{equation}

\begin{equation}
dy = (joyY - CENTER)\cdot k
\end{equation}

kde $k$ je koeficient citlivosti.

\subsection{Komunikácia s riadiacou úlohou}

UI komunikuje s riadiacou úlohou pomocou FreeRTOS front:

\begin{itemize}
    \item \texttt{gUiToControlQueue} -- odosielanie požadovanej polohy/orientácie/kĺbových uhlov
    \item \texttt{gControlToUiQueue} -- prijímanie aktuálneho stavu robota
\end{itemize}

Ak je robot v režime polohy/orientácie:

\begin{equation}
[x,y,z,R,P,Y]
\end{equation}

sa odošlú do inverznej kinematiky.

Ak je robot v režime kĺbového ovládania:

\begin{equation}
[\theta_1,\theta_2,\theta_3,\theta_4,\theta_5,\theta_6]
\end{equation}

sa odošlú priamo regulátoru servomotorov.

\subsection{Chybové hlásenia}

Pri výskyte chyby sa na displeji zobrazí chybová správa.

\begin{lstlisting}[caption={Ukážka chyby}]
CONTROL ERROR
IK FAILED
\end{lstlisting}

alebo

\begin{lstlisting}[caption={Ukážka chyby}]
CONTROL ERROR
OUT OF LIMITS
\end{lstlisting}

Týmto spôsobom používateľ okamžite vidí, že požadovaná poloha nie je dosiahnuteľná alebo prekračuje limity robota.